% Standalone source generated by build.py; edit article.tex and sections/.
\documentclass[11pt,letterpaper]{article}
\usepackage[T1]{fontenc}
\usepackage{lmodern,microtype}
\usepackage[margin=1in,headheight=14pt]{geometry}
\usepackage{amsmath,amssymb,amsthm,mathtools,bm}
\usepackage{mathrsfs}
\usepackage{booktabs,longtable,array}
\usepackage{xcolor}
\definecolor{ink}{HTML}{18384B}
\usepackage[colorlinks=true,linkcolor=ink,citecolor=ink,urlcolor=ink,breaklinks=true]{hyperref}
\usepackage{fancyhdr}
\pagestyle{fancy}\fancyhf{}
\fancyhead[L]{\small Finite-difference determinants}
\fancyhead[R]{\small ProveIt research continuation}
\fancyfoot[C]{\thepage}
\renewcommand{\headrulewidth}{0.3pt}
\setlength{\parskip}{0.3em}
\setlength{\parindent}{1.2em}
\allowdisplaybreaks[1]
\numberwithin{equation}{section}
\newtheorem{theorem}{Theorem}[section]
\newtheorem{lemma}[theorem]{Lemma}
\newtheorem{proposition}[theorem]{Proposition}
\newtheorem{corollary}[theorem]{Corollary}
\theoremstyle{definition}\newtheorem{definition}[theorem]{Definition}
\theoremstyle{remark}\newtheorem{remark}[theorem]{Remark}
\DeclareMathOperator{\Li}{Li}
\DeclareMathOperator{\FP}{FP}
\DeclareMathOperator{\rank}{rank}
\DeclareMathOperator{\cof}{cof}
\DeclareMathOperator{\Rea}{Re}
\newcommand{\C}{\mathbb C}\newcommand{\Q}{\mathbb Q}\newcommand{\N}{\mathbb N}
\newcommand{\Z}{\mathcal Z}\newcommand{\Bell}{\mathcal B}
\newcommand{\dd}{\,\mathrm d}
\newcommand{\ee}{\bm e}
\newcommand{\sym}{\mathrm{sym}}
\newcommand{\pin}{\texttt{73f18ef013d5b74eafc883efb2c15210642df12a}}
\newcommand{\fall}[2]{#1^{\underline{#2}}}
\hypersetup{pdftitle={Finite-Difference Determinants: Single-Zeta Closure, Repeated-Pole Classification, and Stieltjes Jets},pdfauthor={OpenAI-assisted research draft for Vladimir Reshetnikov and the ProveIt project},pdfsubject={Exact identities for harmonic determinant sums, Hurwitz zeta, and polylogarithms}}
\begin{document}
\hypersetup{pageanchor=false}
\begin{titlepage}
\centering
\vspace*{0.8cm}
{\small\scshape Analysis / Polylogarithms\par}
\vspace{0.9cm}
{\huge\bfseries\color{ink} Finite-Difference\\[0.18em] Determinants\par}
\vspace{0.7cm}
{\Large Single-zeta closure, repeated-pole classification,\\[0.25em] and Stieltjes jets\par}
\vspace{0.8cm}
{\large A research continuation for Vladimir Reshetnikov\par}
{and the ProveIt project\par}
\vspace{0.25cm}
{OpenAI-assisted mathematical research draft\par}
{11 October 2026\par}
\vspace{0.8cm}
\begin{minipage}{0.92\textwidth}
\small
\textbf{Abstract.}
We solve a rank-one harmonic-data classification problem raised in the
repository's determinantal harmonic-resonance report. For proper rational
weights on finitely many translated integer lattices, we give an explicit
necessary and sufficient residue test. At nonzero critical decay
$x^{-(2r-1)}$, the criterion is particularly simple: rank one is equivalent
to being a $(2r-2)$nd finite difference of a proper rational function.
Repeated poles and several residue classes are allowed.

This algebraic structure produces an independent-complex-order family of
Vandermonde-square sums. At every particle number its diagonal specialization
is a finite exponential coefficient plus one Hurwitz zeta value. Polarization
allows independent orders and shifts without producing products of zeta
functions. Integer orders yield Gamma-weighted complete-harmonic identities.
A normally convergent largest-index subtraction gives all Stieltjes jets,
with the correct distinction between a hard-cutoff constant and a spectral
finite part. The same rank calculation bounds the depth of rational-weight
cutoff generators; an explicit repeated-pole example is evaluated in
$\Li_{2,2}$ and ordinary polylogarithms.

All infinite-family statements have mathematical proofs. The accompanying
programs record 1,037 exact assertions and 29 numerical diagnostics; the
latter are not interval certificates. Novelty is assessed against a targeted
source audit, not the entire literature. The Gaussian $S_6$ and revised
$S_8$ conjectures are not resolved.
\end{minipage}
\vfill
{\small Repository source anchor: \pin\par}
{\small Sources, inspection limits, and proposed integration are in Section~\ref{sec:audit}.\par}
\end{titlepage}
\pagenumbering{roman}\hypersetup{pageanchor=true}
\begingroup\small\setlength{\parskip}{0pt}\tableofcontents\endgroup
\clearpage
\pagenumbering{arabic}
% BEGIN sections/01-context.tex
\section{The continuation problem and the precise scope}
\label{sec:context}

\subsection{What is already available}
The canonical manuscript connects nested sums, Gamma identities, Hurwitz
zeta jets, and generalized Stieltjes constants. Its preface carefully
separates ordinary mathematical proofs, exact finite algebra, and numerical
evidence \cite{repo:main}. That distinction is particularly relevant here:
a formal rank statement about coefficient matrices does not assert
arithmetic independence of the associated special values.

The preserved \emph{Determinantal Harmonic Resonance} report evaluates a
Gauss--Vandermonde partition function and treats its first logarithmic
resonance. For critical rational weights with simple poles on a common
integer grid, its finite moment matrix contains harmonic dependence in
only one entry. This gives depth-two cutoff generators even when the
number of summation variables is arbitrary \cite{repo:dhr}. Its further
research section asks for a classification with repeated poles, a rank-based
depth theorem, and extensions to several residue classes. Those are the
specific source questions addressed in this article.

Some nearby questions have already moved beyond that report. The current
incoming archive \emph{Polygamma Reconstruction and Resonance} records a
solution of the next two-particle resonance and of higher spatial Mellin
coefficient questions. The current \emph{Harmonic Tails, Shifts and Twists}
archive records the gap-four tail-rank family and further shifted identities
\cite{repo:incoming}. We do not claim those results again. Our principal
spectral deformation is a finite difference of an independently variable
power; it is not the same deformation as a moving denominator parameter
in the earlier Gauss ensemble.

\subsection{The three main advances}
The algebraic classification in Theorems~\ref{thm:classification}
and~\ref{thm:critical} is a necessary and sufficient test, not merely a
construction of examples. At nonzero critical decay, it identifies exactly
when repeated-pole harmonic data remain rank one. Its finite-difference
interpretation explains why arbitrarily high pole multiplicities need not
increase the depth of a cutoff generator.

The spectral identity in Theorem~\ref{thm:single} evaluates a new explicit
family for every particle number and every complex order in its convergence
domain. The polarized version, Theorem~\ref{thm:polarized}, permits independent
orders and independent shifts. Despite this freedom, only single Hurwitz
zeta functions appear: no product of two such functions is needed.

The normal subtraction in Theorem~\ref{thm:subtraction} supplies ordinary
convergent series for every associated Stieltjes jet. Its pole-cancelling
factor is explicit. The hard-cutoff constant and spectral finite part are
then compared in Corollary~\ref{cor:cutoff}; their difference is not discarded.

\subsection{Classical ingredients and novelty qualification}
Finite differences, Cauchy--Binet, Gamma differentiation, and the Hurwitz
shift formula are classical. Rational summability and discrete residues
have an extensive theory, including the work of Chen--Singer and
Arreche--Sitaula \cite{chen-singer,arreche-sitaula}. We do not present the
basic discrete-residue criterion as new. Our use of repeated differences
is tied to a prescribed finite Hankel matrix and its critical decay:
this gives the explicit harmonic-pencil classification and the special-value
identities developed below.

The new statements are research contributions relative to the inspected
repository material. The source inspection and literature checks do not
establish global historical priority. Nor do they establish that the
resulting polylogarithmic representations have minimal arithmetic depth.
The text makes no claim of proof-assistant formalization.

\subsection{Conventions}
We fix an integer $r\ge2$ and write
\begin{equation}\label{eq:d}
 d=2r-2,\qquad \Delta f(x)=f(x+1)-f(x).
\end{equation}
Thus $d$ is even. A sum over $n_1<\cdots<n_r$ is ordered; it has no
additional factor $r!$. Put
\[
 \Delta(\bm n)=\prod_{1\le i<j\le r}(n_j-n_i).
\]
The repeated use of $\Delta$ for a finite-difference operator and a
Vandermonde is distinguished by its argument.

Throughout the spectral sections, $\Rea a>0$ and
$(n+a)^{-s}=\exp(-s\log(n+a))$ uses the principal logarithm. All
parameter derivatives are derivatives of functions on open complex
domains, never derivatives inferred from a formula known only at integer
orders. We write
\[
 (s)_k=s(s+1)\cdots(s+k-1),\qquad
 H_N^{(q)}=\sum_{j=1}^N j^{-q},\qquad H_0^{(q)}=0.
\]
The generalized Stieltjes convention is
\begin{equation}\label{eq:stieltjes}
 \zeta(1+u,a)=\frac1u+
 \sum_{m\ge0}\frac{(-1)^m}{m!}\gamma_m(a)u^m.
\end{equation}
In particular $\gamma_0(a)=-\psi(a)$ and $\gamma_0(1)=\gamma$.
A superscript on $\zeta$ differentiates its first variable unless an
argument derivative is explicitly indicated. We use the standard Hurwitz
and Gamma identities with these conventions \cite{dlmf:hurwitz,dlmf:gamma}.

A matrix pencil is a matrix whose entries are linear forms in independent
formal variables. Its rank means rank over the rational-function field in
those variables. This is not the rank of a matrix obtained by treating
numerical harmonic values as independent numbers.

% END sections/01-context.tex
% BEGIN sections/02-spectral.tex
\section{An all-particle single-zeta determinant identity}
\label{sec:spectral}

\subsection{Weights and convergence}
Define
\begin{equation}\label{eq:weight}
 w_n(s,a)=\frac{\Delta^d(n+a)^{-s}}{d!}
 =\frac1{d!}\sum_{j=0}^{d}(-1)^j\binom dj(n+a+j)^{-s}.
\end{equation}
For each $n$, this is entire in $s$ and holomorphic for $\Rea a>0$.
Its finite-difference form, rather than a generic Gamma ratio, is the
essential feature.

\begin{lemma}[Decay, positivity, and normal convergence]\label{lem:decay}
Uniformly on compact subsets of $\C\times\{\Rea a>0\}$,
\begin{equation}\label{eq:weight-asymp}
 w_n(s,a)=c_d(s)(n+a)^{-s-d}
          +O\bigl((n+1)^{-\Rea s-d-1}\bigr),
 \qquad c_d(s)=\frac{(s)_d}{d!}.
\end{equation}
For $s,a>0$, $w_n(s,a)>0$. The partition function
\begin{equation}\label{eq:partition}
 \Z_r(s,a)=\sum_{0\le n_1<\cdots<n_r}
       \Delta(\bm n)^2\prod_{i=1}^r w_{n_i}(s,a)
\end{equation}
converges absolutely and locally normally for $\Rea s>1$, $\Rea a>0$.
All fixed finite mixed derivatives may be taken termwise there.
\end{lemma}
\begin{proof}
Repeated use of the fundamental theorem of calculus gives
\[
 \Delta^d f(x)=\int_{[0,1]^d}f^{(d)}(x+t_1+\cdots+t_d)
       \dd t_1\cdots\dd t_d.
\]
Since $d$ is even, $f^{(d)}(x)=(s)_d x^{-s-d}$ for $f(x)=x^{-s}$.
Expanding the shifted power to first order proves
\eqref{eq:weight-asymp}, uniformly on the indicated compact sets, including
at zeros of $(s)_d$.
For $\Rea s>0$ the Laplace representation is
\begin{equation}\label{eq:laplace}
 w_n(s,a)=\frac1{d!\Gamma(s)}
 \int_0^\infty t^{s-1}e^{-(n+a)t}(1-e^{-t})^d\dd t.
\end{equation}
It proves strict positivity for positive real parameters.

For ordered nonnegative indices,
$|\Delta(\bm n)|^2\le\prod_i(n_i+1)^{2(i-1)}$.
After dropping the ordering, a majorant for \eqref{eq:partition} is
therefore a product of convergent one-dimensional series: its least
rapidly decaying factor has exponent $-\Rea s$.
Local normal convergence follows for $\Rea s>1$. Differentiation is
justified on slightly larger compact neighborhoods by the same argument,
or by the corresponding logarithmic factors and Cauchy estimates.
\end{proof}

\subsection{The finite endpoint matrix}
For $0\le k\le d$, set
\begin{equation}\label{eq:E}
 E_k(s,a)=-\frac1{d!}\sum_{j=1}^{d}(-1)^j\binom dj
       \sum_{\ell=0}^{j-1}(\ell-j)^k(\ell+a)^{-s}.
\end{equation}
Define
\begin{equation}\label{eq:UV}
 E(s,a)=(E_{i+j}(s,a))_{0\le i,j<r},\qquad
 U_r(s,a)=\det E(s,a),\qquad
 V_r(s,a)=\det(E_{i+j}(s,a))_{0\le i,j<r-1}.
\end{equation}
These are finite expressions, entire in $s$. Their coefficients involve
only finitely many powers $(a+\ell)^{-s}$; no infinite special-function
coordinate is hidden in $U_r$ or $V_r$.

\begin{lemma}[Moment cancellation]\label{lem:moments}
For $\Rea s>1$ and $0\le k\le d$,
\begin{equation}\label{eq:moment}
 \mu_k(s,a):=\sum_{n\ge0}n^k w_n(s,a)
       =E_k(s,a)+\mathbf1_{k=d}\zeta(s,a).
\end{equation}
More generally, for every integer $N\ge0$ and every $s\in\C$,
\begin{align}
 \sum_{n=0}^Nn^k w_n(s,a)
 ={}&E_k(s,a)+\mathbf1_{k=d}\sum_{n=0}^N(n+a)^{-s}
       +R_{k,N}(s,a),\label{eq:finite-moment}\\
 R_{k,N}(s,a)
 ={}&\frac1{d!}\sum_{j=1}^d(-1)^j\binom dj
       \sum_{\ell=N+1}^{N+j}(\ell-j)^k(\ell+a)^{-s}.
       \label{eq:upper-boundary}
\end{align}
Here $0^0=1$ in a zeroth moment.
\end{lemma}
\begin{proof}
The polynomial finite-difference identities are
\begin{equation}\label{eq:annihilator}
 \sum_{j=0}^d(-1)^j\binom dj(x-j)^k
 =\begin{cases}0,&k<d,\\ d!,&k=d.\end{cases}
\end{equation}
Insert \eqref{eq:weight}, shift $\ell=n+j$, and use
\[
 \sum_{\ell=j}^{N+j}F(\ell)
 =\sum_{\ell=0}^{N}F(\ell)
  -\sum_{\ell=0}^{j-1}F(\ell)
  +\sum_{\ell=N+1}^{N+j}F(\ell).
\]
This finite identity remains valid when the subtracted and added ranges
overlap. It proves \eqref{eq:finite-moment} for all $s$.

For the infinite identity, first take $\Rea s>d+1$, so all separately
shifted series converge. Equation~\eqref{eq:annihilator} leaves only
$\zeta(s,a)$ at $k=d$ and the finite lower boundary \eqref{eq:E}.
Both sides of \eqref{eq:moment} are holomorphic on $\Rea s>1$ by
Lemma~\ref{lem:decay}, so the identity theorem proves the full stated
domain. This step avoids an illicit splitting into divergent series.
\end{proof}

\begin{theorem}[Single-zeta closure at every particle number]\label{thm:single}
For $r\ge2$, $\Rea a>0$, and $\Rea s>1$,
\begin{equation}\label{eq:single}
 \boxed{\displaystyle\Z_r(s,a)=U_r(s,a)+V_r(s,a)\zeta(s,a).}
\end{equation}
Consequently \eqref{eq:partition} has a meromorphic continuation to the
whole $s$-plane with at most a simple pole at $s=1$. Its residue is
$V_r(1,a)$. For $\Rea a>0$ this residue is nonzero and equals
\begin{equation}\label{eq:residue-product}
 V_r(1,a)=\prod_{j=0}^{r-2}
 \frac{j!\,\Gamma(a+j)\Gamma(j+1)\Gamma(r-j)}
 {\Gamma(2r-1-j)\Gamma(a+2r-2-j)}.
\end{equation}
For positive real $a$, it is strictly positive.
\end{theorem}
\begin{proof}
Finite Cauchy--Binet applied to the monomial evaluation matrix gives
\[
 \sum_{0\le n_1<\cdots<n_r\le N}\Delta(\bm n)^2\prod_i w_{n_i}
   =\det\left(\sum_{n=0}^N n^{i+j}w_n\right)_{0\le i,j<r}.
\]
Pass to the limit using normal convergence and
Lemma~\ref{lem:moments}. The limiting matrix is
\[
 E(s,a)+\zeta(s,a)\ee\ee^{\mathsf T},\qquad
 \ee=(0,\ldots,0,1)^{\mathsf T}.
\]
The determinant is affine in its bottom-right entry, with coefficient
$V_r$. This proves \eqref{eq:single}. The standard meromorphic
continuation of Hurwitz zeta gives the continuation statement
\cite{dlmf:hurwitz}.

At $s=1$, $w_n(1,a)=\Gamma(n+a)/\Gamma(n+a+d+1)$.
The matrix defining $V_r$ is the moment matrix of an $(r-1)$-particle
positive ensemble for $a>0$; its series converges, since its highest
moment has degree $d-2$. Thus $V_r(1,a)>0$. The product evaluation is
a specialization of the Gauss--Vandermonde calculation in the source
report \cite{repo:dhr}; Appendix~\ref{app:product} gives the finite
factorial-determinant derivation with the present normalization. The product
has neither zeros nor poles for $\Rea a>0$, proving nonvanishing there.
\end{proof}

\subsection{An explicit two-particle identity}
For $r=2$, put $A=a^{-s}$ and $B=(a+1)^{-s}$. Directly from
\eqref{eq:E},
\[
 E_0=\frac{A-B}{2},\qquad E_1=\frac B2,\qquad
 E_2=-A-\frac B2.
\]
Therefore
\begin{equation}\label{eq:r2}
 \boxed{\displaystyle
 \Z_2(s,a)=\frac{a^{-s}-(a+1)^{-s}}2\zeta(s,a)
       -\frac{a^{-2s}}2+\frac{[a(a+1)]^{-s}}4.}
\end{equation}
This identity holds initially for $\Rea s>1$ and gives the continuation
elsewhere. At $a=1$ and integer $m\ge2$,
\begin{equation}\label{eq:r2-integer}
 \Z_2(m,1)=\frac{1-2^{-m}}2\zeta(m)-\frac12+2^{-m-2}.
\end{equation}
For example,
\[
 \Z_2(2,1)=\frac{\pi^2-7}{16},\qquad
 \Z_2(3,1)=\frac{14\zeta(3)-15}{32}.
\]
The three-particle specializations include
\begin{align}
 \Z_3(2,1)&=\frac{928\zeta(2)-1325}{5971968},\label{eq:r3m2}\\
 \Z_3(3,1)&=\frac{333120\zeta(3)-390083}{2579890176}.
 \label{eq:r3m3}
\end{align}
These are exact determinant evaluations, not fitted numerical relations.

\subsection{A specialization anomaly that must not be suppressed}
\label{sec:anomaly}
At $s=0$, every weight in \eqref{eq:weight} vanishes, whereas
\eqref{eq:r2} gives
\begin{equation}\label{eq:anomaly}
 \Z_2^{\mathrm{cont}}(0,a)=-\frac14.
\end{equation}
Similarly, the pointwise weights at $s=-1$ vanish for $r=2$, but
\[
 \Z_2^{\mathrm{cont}}(-1,a)=\frac1{24}.
\]
There is no contradiction. A series that is identically zero at an isolated
parameter need not converge normally in a neighborhood of that parameter.
The continuation is defined from the open half-plane $\Rea s>1$.
Consequently, statements outside that half-plane refer to this continuation,
not automatically to a newly specialized pointwise series. This distinction
will also be retained when discussing finite parts and primitives.

% END sections/02-spectral.tex
% BEGIN sections/03-polarization.tex
\section{Independent spectral orders and independent shifts}
\label{sec:polarization}

\subsection{A mixed determinant with no mixed zeta products}
Let $(s_\nu,a_\nu)$, $1\le\nu\le r$, be independent pairs with
$\Rea s_\nu>1$ and $\Rea a_\nu>0$. The finite-difference order is
still $d=2r-2$ for every weight. Define the averaged symmetrization
\begin{equation}\label{eq:sympartition}
 \Z_r^{\sym}(\bm s,\bm a)
 =\frac1{r!}\sum_{\sigma\in S_r}
   \sum_{0\le n_1<\cdots<n_r}\Delta(\bm n)^2
      \prod_{i=1}^r w_{n_i}(s_{\sigma(i)},a_{\sigma(i)}).
\end{equation}
When all pairs coincide, this is exactly $\Z_r(s,a)$, not $r!\Z_r(s,a)$.

For independent formal variables $t_1,\ldots,t_r$, write
\[
 E(\bm t)=\sum_{\nu=1}^r t_\nu E(s_\nu,a_\nu),\qquad
 C(\bm t)=\det E(\bm t)_{0\le i,j<r-1}.
\]
The latter is a homogeneous polynomial of degree $r-1$.

\begin{theorem}[Polarized single-zeta closure]\label{thm:polarized}
The series \eqref{eq:sympartition} is absolutely and locally normally
convergent in the stated domain, and
\begin{align}
 \Z_r^{\sym}(\bm s,\bm a)
 ={}&\frac1{r!}[t_1\cdots t_r]\det E(\bm t)\notag\\
 &+\frac1{r!}\sum_{\nu=1}^r\zeta(s_\nu,a_\nu)
       \left[\prod_{\kappa\ne\nu}t_\kappa\right]C(\bm t).
       \label{eq:polarized}
\end{align}
The coefficient expressions are finite and entire in all spectral
variables. Thus the joint continuation contains only a linear combination
of single Hurwitz zeta functions. It has no products of zeta functions
and no product poles at intersections of $s_\nu=1$.
\end{theorem}
\begin{proof}
Each term in the finite symmetrization satisfies the product majorant from
Lemma~\ref{lem:decay}, using $\min_\nu\Rea s_\nu>1$ on compact subsets.
Apply Cauchy--Binet to the combined weight
$\sum_\nu t_\nu w_n(s_\nu,a_\nu)$. Its moment matrix is
\[
 E(\bm t)+\left(\sum_\nu t_\nu\zeta(s_\nu,a_\nu)\right)
       \ee\ee^{\mathsf T}.
\]
Its determinant is
\[
 \det E(\bm t)+
 \left(\sum_\nu t_\nu\zeta(s_\nu,a_\nu)\right)C(\bm t).
\]
On the series side, the coefficient of $t_1\cdots t_r$ sums all $r!$
assignments of the labeled weights to the ordered particles. Divide by
$r!$. This proves the formula and its precise normalization. The
meromorphic conclusion follows because the displayed finite coefficients
are entire in the independent spectral variables.
\end{proof}

\begin{remark}[Why independent shifts do not increase the rank]
Each separate shift creates its own Hurwitz zeta coordinate, but every
coordinate enters the same bottom-right moment entry. The number of
special-function coordinates and the rank of the matrix subspace carrying
them are different quantities. This is the mechanism behind the
multiple-residue-class extension in Section~\ref{sec:classification}.
\end{remark}

For $r=2$ and a common shift $a$, set
$A_s=a^{-s}$ and $B_s=(a+1)^{-s}$. Equation~\eqref{eq:polarized} reads
\begin{align}
 \Z_2^{\sym}(s,t;a)
 ={}&\frac{A_s-B_s}{4}\zeta(t,a)
     +\frac{A_t-B_t}{4}\zeta(s,a)\notag\\
 &-\frac{A_sA_t}{2}
     +\frac{A_sB_t+A_tB_s}{8}.
 \label{eq:polarized-r2}
\end{align}
In particular,
\begin{equation}\label{eq:mixed23}
 \Z_2^{\sym}(2,3;1)
   =\frac7{32}\zeta(2)+\frac3{16}\zeta(3)-\frac{29}{64}.
\end{equation}
The unaveraged two-term symmetrization is twice this value.

\subsection{All mixed order derivatives}
For every integer $q\ge0$,
\begin{equation}\label{eq:order-weight}
 \partial_s^q w_n(s,a)
 =\frac{(-1)^q}{d!}\sum_{j=0}^d(-1)^j\binom dj
       \frac{\log^q(n+a+j)}{(n+a+j)^s}.
\end{equation}
Differentiating \eqref{eq:polarized} in any fixed collection of the
independent $s_\nu$ therefore evaluates symmetrized sums with independent
finite-difference logarithmic weights. Their right sides contain only
finite coefficients times single derivatives
$\zeta^{(j)}(s_\nu,a_\nu)$.

This is not a formal differentiation of integer-weight identities.
The independent-variable identity is first established on a complex
product domain with normal convergence. Subsequent integer specializations
of its derivatives are consequently justified. Differentiating in the
shifts as well produces mixed logarithmic and polygamma identities
without changing that analytic foundation.

\subsection{A useful distinction from general determinant sums}
A generic determinant of zeta-valued moments contains products of zeta
functions. Here their absence has an exact reason: the coefficient
matrices of all the zeta coordinates share a one-dimensional image.
No arithmetic cancellation conjecture is being used. Conversely, the
formula does not say that an arbitrary unsymmetrized assignment of distinct
orders has the same reduction. The symmetrization is essential to this
particular determinant argument.

% END sections/03-polarization.tex
% BEGIN sections/04-gamma.tex
\section{Gamma-weighted harmonic identities and differential ladders}
\label{sec:gamma}

\subsection{The complete-harmonic weight at every integer order}
For $q\ge1$ define the local harmonic increment
\begin{equation}\label{eq:local-harmonic}
 h_q(n,a)=\sum_{j=0}^d\frac1{(n+a+j)^q}.
\end{equation}
The complete exponential Bell polynomials are normalized by
\begin{equation}\label{eq:bell}
 \exp\left(\sum_{j\ge1}x_j\frac{u^j}{j!}\right)
 =\sum_{k\ge0}\Bell_k(x_1,\ldots,x_k)\frac{u^k}{k!},
 \qquad \Bell_0=1.
\end{equation}

\begin{theorem}[All-particle Gamma--harmonic evaluation]\label{thm:gamma}
For every integer $m\ge1$,
\begin{equation}\label{eq:harmonic-weight}
 w_n(m,a)=\frac{\Gamma(n+a)}{\Gamma(n+a+d+1)}
 \frac{\Bell_{m-1}(h_1,1!h_2,\ldots,(m-2)!h_{m-1})}{(m-1)!}.
\end{equation}
For $m=1$ the Bell factor is interpreted as one. For $m\ge2$, the
ordered $r$-particle sum of the products of these weights satisfies
\begin{align}
 &\sum_{0\le n_1<\cdots<n_r}\Delta(\bm n)^2
 \prod_{i=1}^r
 \left\{\frac{\Gamma(n_i+a)}{\Gamma(n_i+a+d+1)}
 \frac{\Bell_{m-1}(h_1(n_i,a),\ldots,(m-2)!h_{m-1}(n_i,a))}
 {(m-1)!}\right\}\notag\\
 &\hspace{22mm}=U_r(m,a)+
        \frac{(-1)^m}{(m-1)!}V_r(m,a)\psi^{(m-1)}(a).
 \label{eq:gamma-evaluation}
\end{align}
At positive integral $a$, all the $h_q$ are differences of ordinary
generalized harmonic numbers. At rational $a$, they are rational-grid
harmonic increments.
\end{theorem}
\begin{proof}
The elementary beta partial-fraction identity gives
\[
 w_n(1,a)=\frac1{(n+a)_{d+1}}
        =\frac{\Gamma(n+a)}{\Gamma(n+a+d+1)}.
\]
One can also obtain it by integrating
$t^{n+a-1}(1-t)^d/d!$ on $(0,1)$.
Differentiate \eqref{eq:weight} $m-1$ times in $a$:
\begin{equation}\label{eq:weight-derivative-integer}
 w_n(m,a)=\frac{(-1)^{m-1}}{(m-1)!}
          \partial_a^{m-1}w_n(1,a).
\end{equation}
The logarithmic derivatives of the reciprocal rising factorial are
\[
 \partial_a^q\log w_n(1,a)
     =(-1)^q(q-1)!h_q(n,a).
\]
The Bell formula for an exponential derivative, together with its weighted
homogeneity, proves \eqref{eq:harmonic-weight}. Theorem~\ref{thm:single}
and the classical relation
$\zeta(m,a)=(-1)^m\psi^{(m-1)}(a)/(m-1)!$ prove
\eqref{eq:gamma-evaluation} \cite{dlmf:hurwitz}.
\end{proof}

The first weights are particularly transparent:
\begin{align*}
 w_n(2,a)&=\frac{h_1(n,a)}{(n+a)_{d+1}},\\
 w_n(3,a)&=\frac{h_1(n,a)^2+h_2(n,a)}{2(n+a)_{d+1}},\\
 w_n(4,a)&=\frac{h_1(n,a)^3+3h_1(n,a)h_2(n,a)+2h_3(n,a)}
                  {6(n+a)_{d+1}}.
\end{align*}
Thus the theorem evaluates nonlinear harmonic expressions, not only
partial-fraction weights written in an unfamiliar notation.

For $r=2$, set
\[
 L_n=\frac1{n+1}+\frac1{n+2}+\frac1{n+3},\qquad
 P_n=(n+1)(n+2)(n+3).
\]
A concrete identity is
\begin{equation}\label{eq:concrete-harmonic}
 \boxed{\displaystyle
 \sum_{0\le n<k}\frac{(k-n)^2L_nL_k}{P_nP_k}
       =\frac{\pi^2-7}{16}.}
\end{equation}
Replacing each $L_n$ by
$\bigl(L_n^2+(n+1)^{-2}+(n+2)^{-2}+(n+3)^{-2}\bigr)/2$
gives $(14\zeta(3)-15)/32$. At three particles, the local interval
has five terms and the same construction gives
\eqref{eq:r3m2}--\eqref{eq:r3m3}.

\subsection{Spatial derivatives as exact insertion identities}
For complex $s$ and every integer $q\ge0$,
\begin{equation}\label{eq:shift-weight}
 \partial_a^q w_n(s,a)=(-1)^q(s)_q w_n(s+q,a).
\end{equation}
Consequently, on the original convergence domain,
\begin{align}
 \partial_a^q\Z_r(s,a)
 ={}&(-1)^q\sum_{k_1+\cdots+k_r=q}
    \frac{q!}{k_1!\cdots k_r!}\prod_{i=1}^r(s)_{k_i}\notag\\
 &\qquad\times\sum_{n_1<\cdots<n_r}\Delta(\bm n)^2
       \prod_i w_{n_i}(s+k_i,a).
 \label{eq:insertion-ladder}
\end{align}
All the coefficients and sums are explicit. The right side of
\eqref{eq:single} differentiates to
\begin{equation}\label{eq:spatial-zeta}
 \partial_a^q\Z_r(s,a)=\partial_a^q U_r(s,a)
 +\sum_{j=0}^q\binom qj\partial_a^{q-j}V_r(s,a)
          (-1)^j(s)_j\zeta(s+j,a).
\end{equation}
This uses the standard Hurwitz parameter-derivative law, so again only
single zeta coordinates occur.

For example, let
\[
 \mathcal I_r(s,a)=\sum_{n_1<\cdots<n_r}\Delta(\bm n)^2
        \sum_{i=1}^r w_{n_i}(s+1,a)\prod_{j\ne i}w_{n_j}(s,a).
\]
Then
\begin{equation}\label{eq:first-insertion}
 \mathcal I_r(s,a)=V_r(s,a)\zeta(s+1,a)
   -\frac{\partial_aU_r(s,a)+\partial_aV_r(s,a)\zeta(s,a)}s.
\end{equation}
For $\Rea s>1$, an anchored primitive is therefore
\begin{equation}\label{eq:insertion-primitive}
 \int_{a_0}^{a}\mathcal I_r(s,x)\dd x
       =-\frac{\Z_r(s,a)-\Z_r(s,a_0)}s,
 \qquad \Rea a_0,\Rea a>0,
\end{equation}
with the path lying in the right half-plane. The integral is path
independent because the integrand is holomorphic there. Higher insertion
combinations in \eqref{eq:insertion-ladder} have the analogous anchored
primitives obtained by reducing one derivative order. Thus the harmonic
identities come with differential and antiderivative relations, not just
isolated evaluations.

\subsection{Rational shifts and arithmetic coordinates}
For integer $m$, $U_r(m,a)$ and $V_r(m,a)$ are rational functions of $a$.
At $a\in\Q_{>0}$, the all-particle value therefore lies in the
$\Q$-linear span of $1$ and $\zeta(m,a)$.
For instance, the half-shift identity
$\zeta(s,\tfrac12)=(2^s-1)\zeta(s)$ gives ordinary-zeta reductions
of \eqref{eq:single} at $a=1/2$. General rational shifts may instead
be written with residue-class Dirichlet series or finite root filters.
No claim of arithmetic independence of those coordinates is implied.

For noninteger $s$, the finite exponential factors in $U_r,V_r$ remain
well defined. Their continuation should not be confused with a statement
that the corresponding \emph{cutoff generating function} is always a
finite combination of classical polylogarithms. The latter assertion is
proved below only for rational weights, in particular at integer spectral
orders.

% END sections/04-gamma.tex
% BEGIN sections/05-stieltjes.tex
\section{A normal subtraction and all generalized Stieltjes jets}
\label{sec:stieltjes}

\subsection{The correct escape amplitude}
Write
\[
 \Z_{r,N}(s,a)=\sum_{0\le n_1<\cdots<n_r\le N}
       \Delta(\bm n)^2\prod_iw_{n_i}(s,a),\qquad \Z_{r,-1}=0,
\]
and define the largest-index block
\begin{equation}\label{eq:block}
 B_N(s,a)=\Z_{r,N}(s,a)-\Z_{r,N-1}(s,a).
\end{equation}
It vanishes for $N<r-1$. For $N\ge1$, finite Cauchy--Binet on the
remaining particles gives
\begin{equation}\label{eq:block-det}
 B_N(s,a)=w_N(s,a)
 \det\left(\sum_{n=0}^{N-1}n^{i+j}(N-n)^2w_n(s,a)\right)_{0\le i,j<r-1}.
\end{equation}
The amplitude that must be subtracted is
\begin{equation}\label{eq:amplitude}
 A_r(s,a)=c_d(s)V_r(s,a),\qquad c_d(s)=\frac{(s)_d}{d!},
\end{equation}
not just $V_r(s,a)$ away from $s=1$.

\begin{lemma}[Uniform block remainder]\label{lem:block}
Locally uniformly for $\Rea s>0$, $\Rea a>0$,
\begin{equation}\label{eq:block-asymptotic}
 B_N(s,a)=A_r(s,a)(N+a)^{-s}
       +O\bigl((N+1)^{-\Rea s-1}\bigr).
\end{equation}
In particular the difference and all its fixed parameter derivatives
are locally normally summable in that domain.
\end{lemma}
\begin{proof}
The $(r-1)$-particle ensemble with the same weight has total mass $V_r(s,a)$
for $\Rea s>-1$: its largest required moment has degree $d-2$, and
\eqref{eq:moment} without the zeta term extends to that domain by the
identity theorem. For the present proof it suffices to work on a compact
set with $\Rea s\ge\sigma>0$.

The absolute ensemble with one insertion of $\sum_i n_i$ is uniformly
summable. Indeed, bounding its Vandermonde as in Lemma~\ref{lem:decay},
the largest exponent after the insertion is at most $-\sigma-1$.
For $0\le n_i<N$,
\[
 0\le1-\prod_{i=1}^{r-1}(1-n_i/N)^2
       \le\frac2N\sum_{i=1}^{r-1}n_i.
\]
The sum of the absolute ensemble weights with some $n_i\ge N$ is also
at most $N^{-1}$ times the same first-moment majorant. Therefore
\begin{equation}\label{eq:lower-normalization}
 N^{-d}\frac{B_N(s,a)}{w_N(s,a)}=V_r(s,a)+O(N^{-1}),
\end{equation}
where the quotient notation means the lower-particle polynomial sum,
so no division by a zero weight is needed. This proof uses absolute
majorants and is valid for complex $s,a$ as well.

Combine \eqref{eq:lower-normalization} with
\eqref{eq:weight-asymp} and
$N^d(N+a)^{-d}=1+O(N^{-1})$ to obtain
\eqref{eq:block-asymptotic}. Every bracket is holomorphic. Cauchy
estimates on a slightly larger compact neighborhood give summable
majorants for each fixed derivative; equivalently one can retain the
resulting powers of $\log(N+1)$ in the estimate.
\end{proof}

\begin{theorem}[Ordinary pole-cancelled series]\label{thm:subtraction}
The bracketed series
\begin{equation}\label{eq:T}
 T_r(s,a)=\sum_{N\ge0}\left\{B_N(s,a)-A_r(s,a)(N+a)^{-s}\right\}
\end{equation}
is locally normally convergent for $\Rea s>0$, $\Rea a>0$, and
\begin{equation}\label{eq:Tclosed}
 \boxed{\displaystyle
 T_r(s,a)=U_r(s,a)+(1-c_d(s))V_r(s,a)\zeta(s,a).}
\end{equation}
The apparent pole at $s=1$ is removable. The right side gives an entire
continuation in $s$; the assertion that the series itself is ordinary
and normally convergent is limited to the stated half-plane.
\end{theorem}
\begin{proof}
Normal convergence is Lemma~\ref{lem:block}. For $\Rea s>1$, both
unsubtracted series converge, so Theorem~\ref{thm:single} gives
\eqref{eq:Tclosed} by subtraction. The identity theorem then gives the
whole half-plane $\Rea s>0$. Since $c_d(1)=1$, the unique pole of
Hurwitz zeta cancels. The functions $U_r,V_r,c_d$ are entire in $s$.
\end{proof}

\subsection{An explicit triangular system for every Stieltjes constant}
Let
\begin{equation}\label{eq:jetdefs}
 U_k=\partial_s^kU_r(1,a),\quad
 A_k=\partial_s^kA_r(1,a),\quad
 b_k=\left.\partial_s^k\{(1-c_d(s))V_r(s,a)\}\right|_{s=1}.
\end{equation}
Thus $b_0=0$. Also write $T_k=\partial_s^kT_r(1,a)$.

\begin{theorem}[All-order Stieltjes identities]\label{thm:jets}
For every integer $k\ge0$,
\begin{align}
 T_k
 ={}&\sum_{N\ge0}\left\{\partial_s^kB_N(1,a)
   -\frac1{N+a}\sum_{j=0}^k\binom kj A_j
                   [-\log(N+a)]^{k-j}\right\}\label{eq:Tjets-series}\\
 ={}&U_k+\frac{b_{k+1}}{k+1}
   +\sum_{j=1}^k\binom kj b_j(-1)^{k-j}\gamma_{k-j}(a).
 \label{eq:Tjets-closed}
\end{align}
The series is an ordinary locally normally convergent bracketed series.
For $\Rea a>0$, it can be solved recursively for every $\gamma_m(a)$ because
\begin{equation}\label{eq:b1}
 b_1=-H_dV_r(1,a)\ne0.
\end{equation}
\end{theorem}
\begin{proof}
Differentiate \eqref{eq:T} under its normally convergent sum; Leibniz's
rule yields \eqref{eq:Tjets-series}. Set $s=1+u$ in
\eqref{eq:Tclosed}. Multiply
\[
 \sum_{j\ge1}\frac{b_j}{j!}u^j
 \quad\hbox{by}\quad
 \frac1u+\sum_{m\ge0}\frac{(-1)^m\gamma_m(a)}{m!}u^m.
\]
The pole contributes $b_{k+1}/(k+1)$ to the $k$th derivative. The
remaining product contributes the finite convolution in
\eqref{eq:Tjets-closed}. Since
$c_d'(1)=\sum_{j=1}^d1/j=H_d$, equation~\eqref{eq:b1} follows from
Theorem~\ref{thm:single}.
\end{proof}

In explicit recursive form, putting $k=m+1$ gives
\begin{align}
 \gamma_m(a)=\frac{(-1)^m}{(m+1)b_1}
 \left\{T_{m+1}-U_{m+1}-\frac{b_{m+2}}{m+2}\right.\qquad\notag\\
 \left.{}-\sum_{j=2}^{m+1}\binom{m+1}{j}
       b_j(-1)^{m+1-j}\gamma_{m+1-j}(a)\right\}.
 \label{eq:gamma-invert}
\end{align}
The empty sum at $m=0$ is zero. Every coefficient in this triangular
system is a finite combination of powers and logarithms at $a,a+1,\ldots,a+d-1$
and of ordinary harmonic numbers of length $d$.

The amplitude derivatives have a compact all-order description:
\begin{equation}\label{eq:cd-bell}
 \log c_d(1+u)=\sum_{q\ge1}\frac{(-1)^{q-1}}qH_d^{(q)}u^q,
 \qquad |u|<1.
\end{equation}
Consequently
\[
 c_d^{(k)}(1)=\Bell_k\bigl(H_d,-H_d^{(2)},2!H_d^{(3)},\ldots,
                  (-1)^{k-1}(k-1)!H_d^{(k)}\bigr).
\]
Leibniz's rule with the finite determinants $U_r,V_r$ now generates
all the rows without a new infinite summation problem.

\subsection{Concrete low-order identities}
For $r=2$, $a=1$, put $L=\log2$. The finite coefficient functions are
\[
 U(s)=-\frac12+\frac{2^{-s}}4,\qquad
 V(s)=\frac{1-2^{-s}}2,\qquad c_2(s)=\frac{s(s+1)}2.
\]
Equations~\eqref{eq:Tjets-series}--\eqref{eq:Tjets-closed} give
\begin{align}
 T_0&=-\frac34,\label{eq:T0example}\\
 T_1&=-\frac18-\frac12L-\frac38\gamma,\label{eq:T1example}\\
 T_2&=\frac12L^2-\frac14L-\frac{1+3L}{4}\gamma
          +\frac34\gamma_1.\label{eq:T2example}
\end{align}
For example, $A_0=1/4$ and $A_1=3/8+L/4$, so the second line means
precisely
\begin{equation}\label{eq:T1seriesexample}
 \sum_{N\ge0}\left\{B_N'(1,1)
 -\frac{3/8+L/4-\tfrac14\log(N+1)}{N+1}\right\}
       =-\frac18-\frac12L-\frac38\gamma.
\end{equation}
Here $B_N'$ is computed by differentiating the finite two-particle block,
using \eqref{eq:order-weight}. The cancellation belongs to each displayed
bracket; the two component series must not be split at $s=1$.

\subsection{Hard cutoffs are not spectral finite parts}
\begin{corollary}[The cutoff constant and the amplitude derivative]
\label{cor:cutoff}
For $a>0$, put $U_0=U_r(1,a)$ and $V_0=V_r(1,a)$. Then
\begin{equation}\label{eq:hardcutoff}
 \lim_{N\to\infty}\left\{\Z_{r,N}(1,a)-V_0\log(N+a)\right\}
       =U_0+V_0\{\gamma_0(a)-H_d\}.
\end{equation}
The spectral Laurent constant is
\begin{equation}\label{eq:spectralFP}
 \FP_{s=1}\Z_r(s,a)=U_0+\partial_sV_r(1,a)+V_0\gamma_0(a).
\end{equation}
The spectral constant minus the hard-cutoff constant equals
\begin{equation}\label{eq:FPdifference}
 \partial_sV_r(1,a)+H_dV_0=\partial_sA_r(1,a).
\end{equation}
\end{corollary}
\begin{proof}
The finite sum of the brackets in \eqref{eq:T} at $s=1$ is
\[
 \Z_{r,N}(1,a)-V_0\sum_{n=0}^N(n+a)^{-1}.
\]
It tends to $T_0=U_0-H_dV_0$. The digamma shift identity gives
\[
 \sum_{n=0}^N(n+a)^{-1}=\psi(N+a+1)-\psi(a),
\]
and $\psi(N+a+1)-\log(N+a)\to0$. This proves
\eqref{eq:hardcutoff}. Expanding \eqref{eq:single} at $s=1$ gives
\eqref{eq:spectralFP}, and subtraction proves \eqref{eq:FPdifference}.
\end{proof}

For $r=2,a=1$, the constants are respectively
\begin{equation}\label{eq:cutoff-example}
 \frac\gamma4-\frac34,
 \qquad
 \frac\gamma4-\frac38+\frac{\log2}{4}.
\end{equation}
Their difference is $3/8+\log2/4$. The earlier determinantal report
already warns that residue-amplitude derivatives affect this comparison
\cite{repo:dhr}; the present calculation implements that warning for the
new finite-difference spectral deformation. It is not a claim of an error
in that source.

\subsection{Stieltjes primitives with fixed normalization}
For any $m\ge0$ and $\Rea a,\Rea a_0>0$,
\begin{equation}\label{eq:stieltjes-primitive}
 \int_{a_0}^a\gamma_m(x)\dd x
 =\frac{(-1)^{m+1}}{m+1}
    \left\{\zeta^{(m+1)}(0,a)-\zeta^{(m+1)}(0,a_0)\right\}.
\end{equation}
Indeed, comparing coefficients at $s=0$ in
$\partial_a\zeta(s,a)=-s\zeta(s+1,a)$ proves the derivative of
\eqref{eq:stieltjes-primitive}; the anchor fixes its constant. This
classical Hurwitz-jet primitive is consistent with the canonical
manuscript's conventions \cite{repo:integration,dlmf:hurwitz}.
For $m=0$, Lerch's identity reduces it to
\[
 \int_{a_0}^a\gamma_0(x)\dd x=-\log\Gamma(a)+\log\Gamma(a_0).
\]
Substituting \eqref{eq:gamma-invert} supplies ordinary determinant-series
representatives for the integrands of these primitives. Local normal
convergence permits termwise integration along every compact path in
the right half-plane. Thus the Stieltjes series, Gamma primitive, and
spectral derivative descriptions have one consistent normalization.

% END sections/05-stieltjes.tex
% BEGIN sections/06-classification.tex
\section{Classification of rank-one harmonic matrix pencils}
\label{sec:classification}

\subsection{The rational-weight model, including several lattices}
Let $K\subseteq\C$ be a characteristic-zero field containing the
parameters below. Consider a proper rational function
\begin{equation}\label{eq:rational-weight}
 R(x)=\sum_{\nu=1}^{L}\sum_{j=0}^{J_\nu}
             \sum_{h=1}^{m_\nu}\frac{c_{\nu,j,h}}{(x+a_\nu+j)^h},
 \qquad \Rea a_\nu>0.
\end{equation}
Representatives from distinct lattices are chosen inequivalent modulo
$\mathbb Z$, so that the associated harmonic families are not artificially
duplicated. Zero coefficients and gaps in a lattice are allowed.
For $q\ge1$ put
\begin{equation}\label{eq:shifted-harmonic}
 H_N^{(q)}(a_\nu)=\sum_{n=0}^N(n+a_\nu)^{-q}.
\end{equation}
The finite moment matrix is
\begin{equation}\label{eq:rational-moment-matrix}
 M_N=\left(\sum_{n=0}^Nn^{i+j}R(n)\right)_{0\le i,j<r}.
\end{equation}

\begin{proposition}[Explicit harmonic coefficient matrices]\label{prop:pencil}
There is a matrix $D_N$ with entries rational in $N$ over $K$ such that
\begin{equation}\label{eq:pencil-decomposition}
 M_N=D_N+\sum_{\nu=1}^L\sum_{q=1}^{m_\nu}
                H_N^{(q)}(a_\nu)\,C_{\nu,q}.
\end{equation}
Each $C_{\nu,q}$ is a fixed symmetric Hankel matrix. Writing
$(C_{\nu,q})_{ij}=c_{\nu,q}[i+j]$, its entries are
\begin{equation}\label{eq:Centries}
 c_{\nu,q}[k]=\sum_{h=q}^{m_\nu}\binom{k}{h-q}
       \sum_{j=0}^{J_\nu}c_{\nu,j,h}(-a_\nu-j)^{k-h+q},
       \quad 0\le k\le d.
\end{equation}
A summand is omitted unless $0\le h-q\le k$.
\end{proposition}
\begin{proof}
Expand $x^k=((x+a_\nu+j)-(a_\nu+j))^k$. The coefficient of
$(x+a_\nu+j)^{-q}$ in $x^kR(x)$ is exactly the corresponding summand
in \eqref{eq:Centries}. Summing over $n=0,\ldots,N$ and shifting the
finite harmonic sum changes it by
\[
 -\sum_{\ell=0}^{j-1}(\ell+a_\nu)^{-q}
 +\sum_{\ell=N+1}^{N+j}(\ell+a_\nu)^{-q},
\]
which is rational in $N$ over $K$. The polynomial part of $x^kR(x)$
sums to a polynomial in $N$, for instance by the finite Bernoulli-polynomial
summation formula. All remaining terms therefore belong to $D_N$.
\end{proof}

\begin{definition}
The \emph{harmonic matrix pencil} associated with $R$ and $r$ is
\begin{equation}\label{eq:harmonic-pencil}
 \mathscr C(\bm X)=\sum_{\nu,q}X_{\nu,q}C_{\nu,q},
\end{equation}
where the variables $X_{\nu,q}$ are algebraically independent. Its rank
is the generic matrix rank over $K(\bm X)$.
\end{definition}
This definition is intrinsic after the distinct shift-orbits have been
fixed. It does not posit algebraic independence of the actual functions
$H_N^{(q)}(a_\nu)$ or of their limits.

\subsection{The complete rank-one test}
For $\tau\in K$, write
\[
 v(\tau)=(1,\tau,\ldots,\tau^{r-1})^{\mathsf T},\qquad
 \ee=(0,\ldots,0,1)^{\mathsf T}.
\]

\begin{theorem}[All rank-one harmonic pencils]\label{thm:classification}
The nonzero pencil \eqref{eq:harmonic-pencil} has rank one if and only
if exactly one of the following forms is possible:
\begin{align}
 C_{\nu,q}&=\lambda_{\nu,q}v(\tau)v(\tau)^{\mathsf T}
       &&\text{for a common finite }\tau,\label{eq:geometric-branch}\\
 C_{\nu,q}&=\lambda_{\nu,q}\ee\ee^{\mathsf T}
       &&\text{for all }\nu,q.\label{eq:corner-branch}
\end{align}
In either case some $\lambda_{\nu,q}$ is nonzero. Equivalently, the
finite residue sums \eqref{eq:Centries} obey either
\[
 c_{\nu,q}[k]=\lambda_{\nu,q}\tau^k\quad(0\le k\le d)
\]
with one common $\tau$, or
\[
 c_{\nu,q}[k]=\lambda_{\nu,q}\mathbf1_{k=d}\quad(0\le k\le d).
\]
These are necessary and sufficient algebraic tests for the original
rational weight, including repeated poles and multiple lattices.
\end{theorem}
\begin{proof}
If a symmetric pencil has rank at most one, every coefficient matrix
has rank at most one. Two nonzero symmetric rank-one matrices can be
written $\alpha uu^{\mathsf T}$ and $\beta vv^{\mathsf T}$ over a
field extension. If $u,v$ are not proportional, some pair of coordinates
has $u_iv_j-u_jv_i\ne0$. The corresponding principal minor of
$\alpha uu^{\mathsf T}+t\beta vv^{\mathsf T}$ is
\[
 t\alpha\beta(u_iv_j-u_jv_i)^2\ne0.
\]
This contradicts rank at most one. All coefficient matrices consequently
have one common rank-one direction.

It remains to classify a nonzero rank-one symmetric Hankel matrix.
Write it as $\lambda uu^{\mathsf T}$. If $u_0\ne0$, normalize
$u_0=1$. Equality of the entries $(0,j+1)$ and $(1,j)$ gives
$u_{j+1}=u_1u_j$, hence $u_j=\tau^j$. If $u_0=0$, let $k$ be
the first nonzero index. Were $k<r-1$, the Hankel entries $(k,k)$
and $(k-1,k+1)$ would respectively be nonzero and zero. Thus
$k=r-1$, giving the corner direction. The finite parameter $\tau$
belongs to $K$: it is the ratio of entries $(0,1)$ and $(0,0)$ of
any coefficient matrix whose $(0,0)$ entry is nonzero. The converse
is immediate from either displayed form.
\end{proof}

The test is constructive. If some $C_{\nu,q}(0,0)\ne0$, compute
$\tau=C_{\nu,q}(0,1)/C_{\nu,q}(0,0)$ and verify
\eqref{eq:geometric-branch} in all entries. If no such coefficient
exists, only \eqref{eq:corner-branch} can hold. A zero pencil is
reported separately, not mislabeled as rank one.

\subsection{The critical-decay theorem: repeated finite differences}
Let the residue polynomial for pole order $q$ on lattice $\nu$ be
\begin{equation}\label{eq:Pq}
 P_{\nu,q}(z)=\sum_{j=0}^{J_\nu}c_{\nu,j,q}z^j.
\end{equation}

\begin{lemma}[Corner support and finite-difference divisibility]
\label{lem:corner-divisibility}
The following statements are equivalent, with no critical-decay assumption:
\begin{enumerate}
\item Every $C_{\nu,q}$ is supported in its bottom-right entry.
\item For every $\nu,q$ and $0\le k<d$,
      $\sum_jc_{\nu,j,q}(-a_\nu-j)^k=0$.
\item Every $P_{\nu,q}(z)$ is divisible by $(z-1)^d$ in $K[z]$.
\item There exists a unique proper rational function $F$ on the same
      finite collection of translated lattices such that
      $R=\Delta^dF/d!$.
\end{enumerate}
The harmonic pencil is zero exactly when every $P_{\nu,q}$ is
divisible by $(z-1)^{d+1}$, equivalently when $R$ is a
$(d+1)$st finite difference of a proper rational function.
\end{lemma}
\begin{proof}
Assume corner support. At the highest pole order $q=m_\nu$,
\eqref{eq:Centries} states directly that the moments of the residues
of order $q$ vanish through degree $d-1$. Descend in $q$.
Every already treated term with $h>q$ involves a residue moment of
degree $k-h+q<k<d$ and hence vanishes. The remaining term is exactly
the required order-$q$ moment. This proves the implication from the
first to the second statement. Substitution in \eqref{eq:Centries}
proves the converse.

Shifting by $a_\nu$ does not change the span of polynomials of degree
less than $d$. Ordinary powers of $j$ and the falling factorials
$\fall{j}{k}$ span that same space. Thus the moment conditions are
equivalent to $P_{\nu,q}^{(k)}(1)=0$ for $k<d$, which is precisely
the asserted divisibility.

For a proper rational function $F$ with residue polynomials
$Q_{\nu,q}$, replacing $x$ by $x+1$ multiplies each residue polynomial
by $z$. Therefore $\Delta^dF/d!$ has residue polynomials
$(z-1)^dQ_{\nu,q}/d!$. Dividing $P_{\nu,q}$ explicitly constructs
$F$ and proves the fourth statement. Multiplication by $(z-1)^d$
is injective on polynomial residue data, so this proper rational $F$
is unique. A polynomial kernel is excluded by properness.

For a zero pencil repeat the same descending argument through degree
$d$, rather than $d-1$. This gives divisibility by $(z-1)^{d+1}$
and its stated finite-difference interpretation.
\end{proof}

\begin{theorem}[Critical rational weights: an exact if-and-only-if]
\label{thm:critical}
Suppose $R$ has the form \eqref{eq:rational-weight} and
\begin{equation}\label{eq:critical-decay}
 R(x)=\kappa x^{-(d+1)}+O(x^{-(d+2)}),\qquad\kappa\ne0,
 \qquad d=2r-2.
\end{equation}
Then its harmonic pencil has rank one if and only if
\begin{equation}\label{eq:critical-difference}
 \boxed{\displaystyle R(x)=\frac1{(2r-2)!}\Delta^{2r-2}F(x)}
\end{equation}
for a proper rational function $F$ on the same translated lattices.
Equivalently, every residue polynomial \eqref{eq:Pq} is divisible
by $(z-1)^{2r-2}$. In that event every harmonic coefficient lies in
the single bottom-right entry of the moment matrix.
\end{theorem}
\begin{proof}
The coefficient of $x^{-1}$ in the expansion of $x^kR(x)$ at infinity
is the sum of the coefficients of all its simple partial-fraction
terms. By \eqref{eq:critical-decay}, it is zero for $k<d$ and
$\kappa$ for $k=d$. Proposition~\ref{prop:pencil} therefore gives
\begin{equation}\label{eq:sumC1}
 \sum_{\nu=1}^L C_{\nu,1}=\kappa\ee\ee^{\mathsf T}.
\end{equation}
If the harmonic pencil has rank one, the common-image argument in
Theorem~\ref{thm:classification} forces all its coefficient matrices
to be multiples of this nonzero corner matrix. Apply
Lemma~\ref{lem:corner-divisibility}. Conversely, the finite-difference
condition gives corner support. Equation~\eqref{eq:sumC1} guarantees
that the pencil is nonzero, hence has rank exactly one.
\end{proof}

This answers the source's repeated-pole rank-one question at its
critical decay and supplies a complete finite residue test even away
from that decay. The extension to several residue classes does not
require a new determinant argument: equation~\eqref{eq:sumC1} sums
over those classes, while the finite-difference divisibility is checked
separately on each shift-orbit.

If
\[
 F(x)=\sum_{\nu,j,q}\frac{b_{\nu,j,q}}{(x+a_\nu+j)^q},
 \qquad R=\Delta^dF/d!,
\]
then the corner coefficient is especially simple:
\begin{equation}\label{eq:lambda-seed}
 (C_{\nu,q})_{r-1,r-1}=\sum_j b_{\nu,j,q}.
\end{equation}
Indeed, the $d$th residue moment of a normalized $d$th difference
returns the leading coefficient of that polynomial moment; terms
of higher pole order involve lower residue moments and vanish.
For $F(x)=(x+a)^{-m}$ this coefficient is one at order $m$ and zero
at every other order. This recovers exactly the moment cancellation
of Section~\ref{sec:spectral} at integer spectral orders.

\subsection{Why the nonzero critical coefficient matters}
\label{sec:counterexample}
For $r=2$, consider
\begin{align}
 R(x)&=\frac{2}{(x+1)^2(x+2)(x+3)}\label{eq:counterexample}\\
 &=\frac1{(x+1)^2}-\frac{3}{2(x+1)}
       +\frac2{x+2}-\frac1{2(x+3)}.\notag
\end{align}
It decays as $2x^{-4}$, faster than the critical $x^{-3}$.
Its harmonic coefficient matrices are
\[
 C_1=0,\qquad
 C_2=\begin{pmatrix}1&-1\\-1&1\end{pmatrix}
       =v(-1)v(-1)^{\mathsf T}.
\]
The harmonic pencil has rank one, but $P_2(z)=1$ is not divisible
by $(z-1)^2$. Thus $R$ is not a second difference of a proper
rational function. This is the geometric branch in
Theorem~\ref{thm:classification}, not a counterexample to
Theorem~\ref{thm:critical}.

The example is supplied to prevent an overextension of the present
result. It is not attributed to a false theorem in the inspected
repository. In particular, one must not replace
\eqref{eq:critical-decay} by an unspecified ``at least critical decay''
in the if-and-only-if statement.

\subsection{An exponential-jet version of the test}
The residue tests can be compressed into finite exponential jets:
\begin{equation}\label{eq:jet-pencil}
 \mathcal M_{\nu,q}(t)=e^{-a_\nu t}
     \sum_{h=q}^{m_\nu}\frac{t^{h-q}}{(h-q)!}
             P_{\nu,h}(e^{-t})\pmod{t^{d+1}}.
\end{equation}
Expanding the right side shows that
\[
 \mathcal M_{\nu,q}(t)=\sum_{k=0}^d c_{\nu,q}[k]\frac{t^k}{k!}.
\]
Thus rank one means either
$\mathcal M_{\nu,q}(t)=\lambda_{\nu,q}e^{\tau t}\pmod{t^{d+1}}$
for one common $\tau$, or
$\mathcal M_{\nu,q}(t)=\lambda_{\nu,q}t^d/d!\pmod{t^{d+1}}$.
This formulation gives a short certificate representation without
constructing every matrix entry separately.

% END sections/06-classification.tex
% BEGIN sections/07-polylogs.tex
\section{Rank-controlled polylogarithmic generators}
\label{sec:polylogs}

\subsection{What the harmonic rank does and does not imply}
For the rational weight of Section~\ref{sec:classification}, define
\[
 S_{r,N}(R)=\sum_{0\le n_1<\cdots<n_r\le N}
             \Delta(\bm n)^2\prod_iR(n_i)=\det M_N.
\]
No infinite-sum convergence is needed for this finite identity.
By Proposition~\ref{prop:pencil}, $S_{r,N}$ is a polynomial in
shifted harmonic symbols with rational functions of $N$ as coefficients.

\begin{proposition}[Rank bounds harmonic degree]\label{prop:degree}
If the pencil $\mathscr C(\bm X)$ has generic rank $\rho$, then
\begin{equation}\label{eq:degree-bound}
 \deg_{\bm X}\det(D_N+\mathscr C(\bm X))\le\rho.
\end{equation}
In particular, a rank-one pencil makes $S_{r,N}$ affine in all the
harmonic symbols, even with repeated poles and arbitrarily many particles.
A zero pencil makes $S_{r,N}$ rational in $N$.
\end{proposition}
\begin{proof}
Over $K(N,\bm X)$ a rank-$\rho$ matrix is a product of an $r\times\rho$
and a $\rho\times r$ matrix. The determinant
$\det(D_N+t\mathscr C(\bm X))$ therefore has degree at most $\rho$
in $t$, by multilinearity or by the vanishing of minors of order
larger than $\rho$. The coefficient of $t^k$ is exactly the
homogeneous degree-$k$ part in $\bm X$. This proves the assertion
without any assumption that $D_N$ is invertible.
\end{proof}

The degree can be strictly smaller than $\rho$ because of additional
cancellations involving $D_N$. Neither $\rho$ nor the polynomial degree
proves minimal depth of evaluated periods. These are finite structural
upper bounds.

\subsection{A constructive depth bound on rational grids}
For positive integers $p_1,\ldots,p_\ell$, our multiple-polylogarithm
convention is
\begin{equation}\label{eq:MPL}
 \Li_{p_1,\ldots,p_\ell}(z_1,\ldots,z_\ell)
 =\sum_{n_1>\cdots>n_\ell\ge1}
   \frac{z_1^{n_1}\cdots z_\ell^{n_\ell}}
        {n_1^{p_1}\cdots n_\ell^{p_\ell}}.
\end{equation}
Its depth is $\ell$. In the disk, all identities below use the
common power-series germ; rational Laurent prefactors have removable
singularities only after the full expression is assembled.

\begin{theorem}[Harmonic rank controls cutoff-generator depth]
\label{thm:depth}
Assume every $a_\nu$ is positive rational, and choose a common denominator
$D$ for the pole locations. If the harmonic pencil has rank $\rho$, then
\begin{equation}\label{eq:generator-general}
 G_R(z)=\sum_{N\ge0}S_{r,N}(R)z^{N+1},\qquad |z|<1,
\end{equation}
can be expressed by finitely many cyclotomic multiple polylogarithms
of level dividing $D$ and depth at most $\rho+1$, rational functions
of $z^{1/D}$, and Euler derivatives of such expressions.
Every positive Euler primitive, normalized to vanish at $z=0$, obeys
the same depth bound.

For integer-grid poles, no fractional power or root-of-unity filter
is needed. In particular, the rank-one families in
Theorems~\ref{thm:classification}--\ref{thm:critical} have depth at
most two. A zero pencil has depth at most one.
\end{theorem}
\begin{proof}
Proposition~\ref{prop:degree} reduces each coefficient to a finite sum
of terms
\begin{equation}\label{eq:coefficient-type}
 Q(N)\prod_{j=1}^{\ell}H_N^{(q_j)}(a_{\nu_j}),\qquad \ell\le\rho,
\end{equation}
where $Q$ is rational and has rational-grid poles. All such sequences
have at most polynomial-logarithmic growth, so their generating series
converge normally in the disk.

First consider integral shifts. Expand the product of finite harmonic
sums by partitioning the index set according to equality and strict
ordering. This is the finite quasi-shuffle argument; it yields nested
sums of depth at most $\ell$. At an equal index, different shifted
rational factors are reduced by ordinary partial fractions. Integer
shifts of a cutoff can be removed by subtracting a finite boundary strip.
Removing one such strip either changes a rational coefficient or reduces
the number of nested free indices; it cannot increase depth.

Next decompose $Q(N)$ into a polynomial and partial fractions. For an
unshifted denominator, the outer summation adds at most one strict
index to the nested harmonic sum. Its diagonal contribution has smaller
depth. For example,
\begin{align}
 \sum_{M\ge1}\frac{z^M H_M^{(q)}}{M^p}
    &=\Li_{p,q}(z,1)+\Li_{p+q}(z),\label{eq:basic-nesting}\\
 \sum_{M\ge1}z^M H_M^{(q)}&=\frac{\Li_q(z)}{1-z}.
 \label{eq:basic-harmonic-gf}
\end{align}
Shifting an outer denominator replaces $M$ by $M+j$, followed by the
same finite boundary-strip procedure. Polynomial factors in $M$ are
implemented by the Euler operator $z\partial_z$. These operations
prove the bound $\ell+1$ on an integer grid.

For rational shifts write $a=A/D$ after moving the representative
into $0<a\le1$ and accounting for finitely many initial terms.
Then
\[
 H_N^{(q)}(A/D)=D^q
    \sum_{\substack{1\le m\le DN+A\\m\equiv A\ (D)}}m^{-q}.
\]
Resolve a congruence by the exact filter
\[
 \mathbf1_{m\equiv A\ (D)}
     =\frac1D\sum_{j=0}^{D-1}\omega^{j(m-A)},
 \qquad \omega=e^{2\pi i/D}.
\]
After rescaling indices, strict inequalities between the original
integer indices differ from strict inequalities between the new
residue-class indices only on equal-block diagonals. Those diagonals
are reduced by partial fractions and have fewer free indices. Induction
on their number reduces to strictly nested sums of the form
\eqref{eq:MPL}, with roots of unity in the arguments and $z^{1/D}$
in the outer argument. Depth does not increase. This also explains
why different residue classes can be handled together.

Finally, a positive Euler primitive divides the coefficient of
$z^{N+1}$ by a positive power of $N+1$. It changes only the rational
factor $Q(N)$ in \eqref{eq:coefficient-type}, leaving the maximum
number of harmonic factors unchanged. The same reduction applies.
\end{proof}

This is a representation theorem, not an irreducibility theorem. It
also does not assert that every endpoint reduces to ordinary zeta values
and their products. Repeated poles introduce inner indices larger than
one, for which the endpoint reduction problem is different from the
simple-pole family in the earlier report.

\subsection{An explicit repeated-pole \texorpdfstring{$\Li_{2,2}$}{Li(2,2)} identity}
Take $r=2$, $a=1$, $s=2$, and write
\[
 S_N=\Z_{2,N}(2,1),\qquad M=N+1.
\]
The weight has repeated poles,
\[
 w_n(2,1)=\frac12\left\{\frac1{(n+1)^2}
                    -\frac2{(n+2)^2}+\frac1{(n+3)^2}\right\}.
\]
A finite calculation from \eqref{eq:finite-moment} gives
\begin{equation}\label{eq:cutoff-PQ}
 S_{M-1}=P(M)+Q(M)H_M^{(2)},\qquad M\ge1,
\end{equation}
where
\begin{align}
 P(M)={}&-\frac7{16}-\frac3{4(M+2)}+\frac7{4(M+2)^2}
                   +\frac7{8(M+1)^2}-\frac1{2(M+1)^4},\label{eq:Pexplicit}\\
 Q(M)={}&\frac38-\frac1{2(M+1)^2}+\frac1{2(M+2)^2}.
 \label{eq:Qexplicit}
\end{align}
For example, $M=1$ gives zero, as required because there are not yet
two distinct indices.

\begin{theorem}[Closed repeated-pole cutoff generator]\label{thm:explicitGF}
For $|z|<1$,
\begin{align}
 \sum_{M\ge1}S_{M-1}z^M
 ={}&\frac{1-z}{2z^2}\Li_{2,2}(z,1)
  +\frac{10-7z}{8z^2(1-z)}\Li_2(z)
  +\frac{4z-7}{4z^2}\Li_1(z)\notag\\
 &-\frac{\Li_4(z)}{2z}+\frac1{2z}-\frac7{16(1-z)}.
 \label{eq:explicitGF}
\end{align}
The apparent singularity at $z=0$ is removable, and the resulting
function begins at order $z^2$.
\end{theorem}
\begin{proof}
For completeness, the three finite moments, with $H=H_M^{(2)}$, are
\begin{align*}
 \mu_{0,M-1}&=\frac{M(3M^3+18M^2+39M+28)}{8(M+1)^2(M+2)^2},\\
 \mu_{1,M-1}&=\frac{M(M-1)(M^2+7M+8)}{8(M+1)^2(M+2)^2},\\
 \mu_{2,M-1}&=H-
  \frac{9M^4+70M^3+133M^2+76M}{8(M+1)^2(M+2)^2}.
\end{align*}
The determinant $\mu_0\mu_2-\mu_1^2$ gives
\eqref{eq:cutoff-PQ}--\eqref{eq:Qexplicit} by rational arithmetic.

The shifted harmonic sums needed to sum that expression are
\begin{align*}
 \sum_{M\ge1}\frac{z^M H_M^{(2)}}{(M+1)^2}
     &=z^{-1}\Li_{2,2}(z,1),\\
 \sum_{M\ge1}\frac{z^M H_M^{(2)}}{(M+2)^2}
     &=z^{-2}\left\{\Li_{2,2}(z,1)
          -(1+z)\Li_2(z)-2(1-z)\Li_1(z)+3z\right\}.
\end{align*}
To verify the second line, put $L=M+2$ and remove the missing term
from $H_{L-1}^{(2)}$. The remaining rational factor is reduced using
\[
 \frac1{L^2(L-1)^2}
 =\frac1{L^2}+\frac1{(L-1)^2}+\frac2L-\frac2{L-1}.
\]
The plain shifted sums are
\[
 \sum_{M\ge1}\frac{z^M}{(M+j)^p}
  =z^{-j}\left\{\Li_p(z)-\sum_{\ell=1}^{j}\frac{z^\ell}{\ell^p}\right\}.
\]
Insert these identities and \eqref{eq:basic-harmonic-gf} into
\eqref{eq:cutoff-PQ}. Collecting terms proves \eqref{eq:explicitGF}.
All steps are absolutely convergent in the disk. The original series
proves removability and its initial vanishing; no separate branch
assignments at the apparent poles are made.
\end{proof}

As $z\uparrow1$, Abel's theorem and the convergent sequence $S_N$ give
\begin{equation}\label{eq:AbelGF}
 \lim_{z\uparrow1}(1-z)\sum_{M\ge1}S_{M-1}z^M
     =\frac38\zeta(2)-\frac7{16}=\frac{\pi^2-7}{16}.
\end{equation}
This independently visible endpoint agrees with
\eqref{eq:concrete-harmonic}. The depth-two coordinate in the cutoff
generator disappears from this particular Abel limit because its
coefficient has an additional factor $1-z$. This is an exact cancellation,
not a general theorem of depth-two endpoint reducibility.

\subsection{Constructing new examples without a search}
Choose any proper rational seed $F$ with poles on positive rational
translated grids and set $R=\Delta^{2r-2}F/(2r-2)!$.
Lemma~\ref{lem:corner-divisibility} proves that its cutoff partition has
an affine harmonic expression and hence a depth-two generator at most.
This works for seeds with several pole orders and several residue classes.
If the seed is itself a first difference, the harmonic pencil vanishes
and the generator has depth at most one.

For example, any seed of the form
\[
 F(x)=\frac{u}{x+1/3}+\frac{v}{(x+1/2)^2}
                         +\frac{w}{(x+2/3)^3}
\]
produces an all-particle rank-at-most-one family on level six. The corner
entries are $u,v,w$ in their respective harmonic coordinates by
\eqref{eq:lambda-seed}. When $u\ne0$, the decay is nonzero critical
and Theorem~\ref{thm:critical} applies. The theorem supplies a finite
construction and a complete rank certificate; it does not require a
numerical integer-relation search.

% END sections/07-polylogs.tex
% BEGIN sections/08-audit.tex
\section{Source audit, corrections, and verification status}
\label{sec:audit}

\subsection{The source anchor and the inspection actually performed}
The repository was anchored at commit
\begin{center}\pin,\end{center}
whose timestamp is 11 October 2026 at 18:22:01 UTC. Both directories in
the research request were inspected at that pin. The source record in
\texttt{provenance/sources.json} lists the relevant paths and the level of
inspection, rather than implying a full audit of the repository.

The targeted reading included the canonical main file and preface, the
research-status chapter, the opening 280 lines of the integration chapter,
and the adjacent-tail and research-question sections of the differential
polygamma continuation. The determinantal harmonic-resonance report was
read for its defining ensemble, Gauss product, first-resonance context,
stated simple-pole harmonic-rank result, and further research questions.
Its requested repeated-pole classification and rational-grid extension
are the direct targets of Section~\ref{sec:classification}.

The incoming directory contained six archives at the anchor. The member
list and selected status/readme material of
\emph{Polygamma Reconstruction and Resonance} and the readme of
\emph{Harmonic Tails, Shifts and Twists} were inspected. A combined retrieval
also exposed readme material from \emph{Harmonic Tail Gamma Continuation},
but it was not treated as a full interior audit. The Beta Crossing,
Quadratic Spectrum, and Same-Direction Complex Orders packages were not
read in full in this continuation. No error or absence of a theorem is
attributed to their unread contents. This targeted comparison is enough
to identify the specific source question being answered, but not to prove
an exhaustive novelty claim.

\subsection{Proposed status updates}
The source question about rank-one harmonic pencils should now be split
into a proved classification and remaining arithmetic-depth questions.
Theorem~\ref{thm:classification} supplies the full residue test in the
specified rational-weight model. Theorem~\ref{thm:critical} gives the
particularly short critical-decay characterization, allowing repeated
poles and multiple rational residue classes. The source's simple-pole
construction is thereby extended, not invalidated.

The source question about harmonic rank controlling depth receives the
constructive upper bound in Theorem~\ref{thm:depth}. It should not be
rewritten as a theorem about \emph{minimal} depth. The possible strict
inequality between harmonic rank, determinant degree, and arithmetic depth
remains an important separate issue.

The next-two-particle-resonance question in the older report should be
cross-linked to the current incoming reconstruction/resonance package,
subject to the repository's usual integration review. The present
finite-difference deformation should be listed separately: its simple
pole and its coefficient $c_d(s)$ are not the spectral excess and pole
structure of every higher Gauss resonance.

\subsection{Corrections and safeguards established here}
No newly false analytic theorem was found in the portions of the existing
material inspected for this article. The following are precise
integration safeguards, with explicit examples rather than allegations
about unread sources.

First, the critical-decay characterization requires a \emph{nonzero}
coefficient $\kappa$ in \eqref{eq:critical-decay}. The faster-decaying
weight \eqref{eq:counterexample} has a rank-one geometric pencil and is
not a $d$th finite difference. Replacing exact critical decay by merely
a big-$O$ condition would make the new theorem false.

Second, the escape amplitude is $A_r=c_dV_r$. A counterterm using $V_r$
for all $s$ gives the wrong derivative rows. In particular, the pole-times-
next-coefficient term $b_{k+1}/(k+1)$ in
\eqref{eq:Tjets-closed} must be retained.

Third, the hard-cutoff constant and the spectral Laurent constant differ
by $\partial_sA_r(1,a)$. Equation~\eqref{eq:cutoff-example} gives a
nonzero elementary discrepancy. The earlier source already recognizes
this distinction; our formulas make it explicit for this family.

Fourth, meromorphic continuation does not commute with an isolated
pointwise specialization of a nonnormally convergent family. At $s=0$
the weights vanish, but the continued value is $-1/4$ for two particles.
This also limits the interpretation of identities at nonpositive orders.

Finally, finite harmonic rank and polylogarithmic representation depth
are not arithmetic independence. An explicit $\Li_{2,2}$ term is not
proof that depth two is necessary. Conversely, the simple-pole report's
special ordinary-zeta endpoint closure should not be extrapolated to
all repeated-pole cases without a separate reduction argument.

\subsection{What the supplied verification proves}
The exact suite \texttt{verification/verify\_exact.py} performs
\textbf{1,037 assertions} in rational or symbolic arithmetic. The
machine-readable record is \texttt{results/exact\_checks.json}.
It includes 576 finite endpoint moment identities, 18 direct ordered-sum
Cauchy--Binet comparisons, 18 residue-product checks, 12 single-zeta
matrix identities, independent-order/shift polarization tests, 144
Gamma--harmonic Bell identities, finite-difference residue inversions,
zero-pencil checks, the geometric counterexample, and 20 deliberately
perturbed weights rejected by the rank-one test. It also checks the
explicit cutoff formula, 44 generating-function coefficients, and the
Stieltjes/Bell coefficient rules.

The independent numerical suite records \textbf{29 diagnostics} in
\texttt{results/numerical\_checks.json}. Nine compare determinants of
Mellin-integral moments with the finite-endpoint/zeta expression, including
noninteger and complex orders and shifts. Two test mixed-order/shift
Mellin identities. Four compare direct finite sums with finite boundary
formulas. Four test hard-cutoff constants through finite formulas at
$N=10^{20}$; six test differentiated Stieltjes brackets at $N=10^{18}$.
Three test the explicit polylogarithmic generator, and one tests the
anchored Gamma primitive.

The standard working precision was 60 decimal digits, with increased
precision for cutoff cancellations. The Mellin-determinant relative
residuals in the recorded run are below $6.1\times10^{-57}$.
The largest relative discrepancy across \emph{all} diagnostics is about
$8.83\times10^{-13}$, occurring in a high derivative of a finite-cutoff
approximation. It would therefore be incorrect to describe the whole
suite as agreement to 60 digits. Each row records its actual discrepancy,
tolerance, and cutoff where applicable. These are floating-point
diagnostics, not certified interval enclosures.

The infinite-family conclusions rest on the proofs in this article.
Finite exact checks test implementation and normalization; they do not
replace the general algebraic or analytic arguments. A successful PDF
build likewise verifies document compilation, not mathematical truth.

\subsection{Integration boundary}
The package is self-contained and includes standalone and modular TeX,
the compiled PDF, verification programs and result records, a build
script, a source-scope record, an explicit claim-status file, and checksum
verification. The suggested repository destination is recorded in
\path{integration/INTEGRATION.md}. No repository files were changed,
no archive was retired, and no commit or pull request was created.
The manuscript's Gaussian $S_6$ and revised $S_8$ proof statuses are
unchanged \cite{repo:discovery}.

% END sections/08-audit.tex
% BEGIN sections/09-questions.tex
\section{Further research questions}
\label{sec:questions}

\subsection{Higher-rank harmonic pencils and their moduli}
The rank-one case has the two exact forms
\eqref{eq:geometric-branch}--\eqref{eq:corner-branch}. Classify the
rational residue arrays for which the generic harmonic pencil has rank
at most two, then at most $\rho$. Hankel structure suggests a connection
with finite exponential-polynomial moment sequences and confluent
Vandermonde decompositions. The difficulty is the requirement that
\emph{all} coefficient matrices lie in one low-rank linear space, not
merely that each matrix separately has low rank. An effective necessary
and sufficient classification would extend the present finite residue
certificate beyond rank one.

\subsection{When does determinant degree equal harmonic rank?}
For a given rational weight, let $\delta$ be the total harmonic degree
of $\det(D_N+\mathscr C(\bm X))$. We proved $\delta\le\rho$.
Determine exact conditions for strict inequality, retaining the dependence
of $D_N$ on the same residue data as the pencil. A useful target is a
criterion in minors and discrete residues that avoids expanding the full
determinant. This would sharpen the formal polylogarithmic depth bound
from $\rho+1$ to $\delta+1$ before any functional reductions are used.

\subsection{Minimal depth in a specified functional relation system}
The generator \eqref{eq:explicitGF} contains $\Li_{2,2}(z,1)$, but
its coefficient could interact with functional equations in a broader
alphabet. Determine whether this generator reduces to classical
polylogarithms in a precisely declared differential or symbol calculus.
A nonzero obstruction would concern that formal system, not numerical
period independence. The problem is especially concrete because all
rational coefficients are explicit and the endpoint cancellation is
already known.

\subsection{Closed all-rank formulas for the finite exponential coefficients}
The determinants $U_r,V_r$ are finite at every rank, and $V_r(1,a)$
has the Gamma product \eqref{eq:residue-product}. Find short recurrences
or product formulas for $U_r(m,a)$ and $V_r(m,a)$ at fixed integer
$m\ge2$, or for their jets at $s=1$. The residue-product proof suggests
using finite discrete orthogonal polynomials at $m=1$ followed by
structured confluent perturbations. A successful result should improve
symbolic expression growth without introducing hidden infinite coordinates.

\subsection{Unsymmetrized mixed-order sums}
The polarized determinant evaluates the symmetrized sum with independent
orders and shifts. What additional coordinates are required for an
individual ordered assignment? For two particles, the sum and difference
under exchanging the two labels are natural first targets. The sum is
already evaluated by \eqref{eq:polarized-r2}; the antisymmetric component
may connect to multiple Hurwitz zeta or Lerch functions. One should seek
an identity on a genuine complex-parameter domain before differentiating
at integer values.

\subsection{Noninteger-order cutoff generators}
The spectral identity holds for arbitrary complex orders, but the rational
partial-fraction proof of polylogarithmic cutoff closure uses integer
orders. Determine a natural finite function class for the cutoff generator
at noninteger $s$. A likely candidate involves multiple Hurwitz--Lerch
functions with independently shifted denominators, but merely naming such
a function is not a reduction. The useful target is a finite transform
or differential system that specializes to Theorem~\ref{thm:explicitGF}
and preserves its initial conditions.

\subsection{More general annihilating stencils}
Replace the consecutive finite difference by
$\sum_j\alpha_j f(x+\lambda_j)$, with distinct rational shifts and
$\sum_j\alpha_j\lambda_j^k=0$ for $k<d$. Determine which stencils
still confine the moment special functions to one matrix direction.
When the shifts are integers, residue-polynomial divisibility is the
right language; on nonuniform grids, endpoint terms may produce several
Lerch coordinates. A classification should distinguish a genuinely
rank-one common image from a mere vanishing of the first asymptotic
coefficients.

\subsection{Higher Gauss resonances remain a distinct problem}
For the earlier Gauss ensemble, the resonances
$b-a-c=2r-2-M$ can have several divergent moment entries and overlapping
escape regimes. Our spectral deformation has only the single Hurwitz
pole and therefore does not resolve that general problem. Extend the
largest-index proof with a canonical subtraction of multiple escape
regions, and compare its cutoff polynomial with the Laurent coefficients
of the known Gamma product. The next two-particle case has a separate
incoming treatment; the remaining all-rank statement should retain that
provenance and its more complicated pole structure.

\subsection{Endpoint Euler primitives and weight drops}
Theorem~\ref{thm:depth} bounds the depth of every normalized Euler
primitive inside the disk. Determine which absolutely convergent endpoint
combinations of the integer-order finite-difference family reduce to
ordinary zeta values or products, and which require genuine multiple
zeta coordinates in a stated relation system. The cancellation in
\eqref{eq:AbelGF} is one exact example, not evidence for an unrestricted
closure claim. A systematic endpoint formula would connect the present
cutoff identities with the repository's harmonic-tail and cyclotomic
reduction programs.

\subsection{Formal verification with a small analytic interface}
The finite residue matrices, divisibility tests, polynomial annihilation,
and Cauchy--Binet normalization are suitable exact formalization targets.
They can be separated from three analytic interfaces: local normal
convergence of the spectral ensemble, the uniform largest-index remainder,
and the Hurwitz Laurent expansion. A proof-assistant implementation of
that separation would make subsequent repeated-pole identity generation
substantially more reliable. The supplied scripts are executable
certificates and diagnostics, not such a formalization.

% END sections/09-questions.tex
\appendix
% BEGIN sections/10-residue-product.tex
\section{Derivation of the Gamma product for the residue}
\label{app:product}

The following determinant is a classical finite algebraic ingredient of
the source's Gauss--Vandermonde proof. We include its derivation so the
normalization in \eqref{eq:residue-product} can be checked independently.

\begin{lemma}[A rational factorial determinant]\label{lem:factorial-det}
For indeterminates $x_0,\ldots,x_{L-1},t$,
\begin{equation}\label{eq:factorial-det}
 \det_{0\le i,j<L}\frac{(x_i)_j}{(x_i+t)_j}
 =\frac{\displaystyle\prod_{j=0}^{L-1}(t)_j
             \prod_{0\le i<k<L}(x_k-x_i)}
        {\displaystyle\prod_{i=0}^{L-1}(x_i+t)_{L-1}}.
\end{equation}
It is an identity of rational functions and hence holds wherever the
specialized expressions are defined, with removable values interpreted
by cancellation.
\end{lemma}
\begin{proof}
Multiply row $i$ by $(x_i+t)_{L-1}$. Column $j$ becomes the polynomial
\[
 p_j(x)=(x)_j(x+t+j)_{L-1-j},
\]
of degree at most $L-1$. Its evaluation determinant is the Vandermonde
in the $x_i$ times the determinant of the polynomial coefficients.
Evaluate at $x_i=-i$. The matrix is lower triangular, since $(-i)_j=0$
for $j>i$, and its diagonal is
$(-1)^ii!(t)_{L-1-i}$. Dividing its determinant by the Vandermonde
at $0,-1,\ldots,-(L-1)$ leaves $\prod_{j=0}^{L-1}(t)_j$.
Dividing out the original row factors proves the identity.
\end{proof}

Set $L=r-1$ and $d=2r-2$. Since the two polynomial bases
$\fall ni$ and $(n+a)_j$ are monic in their respective degrees,
Cauchy--Binet gives
\begin{equation}\label{eq:residue-mixed-moment}
 V_r(1,a)=\det_{0\le i,j<L}
 \sum_{n\ge0}\fall ni(n+a)_j
             \frac{\Gamma(n+a)}{\Gamma(n+a+d+1)}.
\end{equation}
Terms $n<i$ vanish. Shift $n=i+k$ and apply Gauss summation
\cite{dlmf:gauss}. Because $i+j\le2r-4=d-2$, its excess
$d-i-j$ is positive. The resulting entry is
\begin{equation}\label{eq:residue-entry}
 M_{ij}=\frac{\Gamma(i+1)\Gamma(a+i+j)\Gamma(d-i-j)}
                    {\Gamma(d+1-j)\Gamma(a+d)}.
\end{equation}
Factor the row terms $\Gamma(i+1)\Gamma(a+i)\Gamma(d-i)$,
the column terms $(-1)^j/\Gamma(d+1-j)$, and one factor
$1/\Gamma(a+d)$ per row. The remaining determinant is
\[
 \det_{0\le i,j<L}\frac{(a+i)_j}{(1-d+i)_j}.
\]
Apply Lemma~\ref{lem:factorial-det} with
$x_i=a+i$ and $t=1-d-a$. None of the denominator rising factorials
vanishes in the required range of indices. Its Vandermonde is
$\prod_{k=0}^{L-1}k!$. Substitution gives the intermediate product
\begin{align}
 V_r(1,a)={}&
 \frac{\prod_{i=0}^{L-1}\Gamma(i+1)\Gamma(a+i)\Gamma(d-i)}
      {\Gamma(a+d)^L}
 \prod_{j=0}^{L-1}\frac{(-1)^j}{\Gamma(d+1-j)}\notag\\
 &\times
 \frac{\prod_{j=0}^{L-1}(1-d-a)_j\prod_{k=0}^{L-1}k!}
      {\prod_{i=0}^{L-1}(1-d+i)_{L-1}}.
 \label{eq:residue-intermediate}
\end{align}
Using
\[
 (1-d-a)_j=(-1)^j\frac{\Gamma(a+d)}{\Gamma(a+d-j)},\qquad
 (1-d+i)_{L-1}=(-1)^{L-1}
       \frac{\Gamma(d-i)}{\Gamma(d-i-L+1)},
\]
and $d-L+1=r$, the signs cancel and
\eqref{eq:residue-intermediate} becomes
\[
 \prod_{j=0}^{r-2}
 \frac{j!\,\Gamma(a+j)\Gamma(j+1)\Gamma(r-j)}
       {\Gamma(2r-1-j)\Gamma(a+2r-2-j)}.
\]
This proves \eqref{eq:residue-product} for $\Rea a>0$. In particular,
\[
 V_2(1,a)=\frac1{2a(a+1)},\qquad
 V_3(1,a)=\frac1{72a(a+1)^2(a+2)^2(a+3)}.
\]
No additional $L!$ is introduced, since the defining ensemble uses
strictly ordered indices throughout.

% END sections/10-residue-product.tex
% BEGIN sections/11-algorithms.tex
\section{Finite algorithms and replayable certificates}
\label{app:algorithms}

\subsection{Computing the single-zeta answer}
For given $r$, order $s$, and shift $a$, compute the $2r-1$ finite
endpoint expressions \eqref{eq:E}. Form their $r\times r$ Hankel matrix
and its upper-left $(r-1)\times(r-1)$ minor. Their determinants are
$U_r,V_r$. The result is then \eqref{eq:single}.
At integer $s=m$, all entries are rational functions of $a$; at rational
$a$, exact rational arithmetic suffices to produce the two coefficients.
The program \texttt{verification/core.py} implements this construction
without evaluating any zeta function.

For a finite cutoff $N$, add the finite harmonic term and the upper
boundary \eqref{eq:upper-boundary} before taking the determinant.
This produces an exact formula whose number of endpoint terms depends
on $r$, not on $N$. The large numerical cutoffs in the verification
record were handled by this formula, not by enumerating $10^{20}$
summands.

\subsection{Classifying a rational weight}
A complete certificate consists of its partial-fraction residues, the
matrices \eqref{eq:Centries}, and one of the following outputs:
zero pencil; corner pencil; geometric pencil with the exact parameter
$\tau$; or a nonzero minor of a linear combination disproving rank one.
At critical decay, the residue-polynomial quotients
\[
 Q_{\nu,q}(z)=d!\,P_{\nu,q}(z)/(z-1)^d
\]
are particularly compact positive certificates. Their coefficients
reconstruct the unique proper rational seed $F$.

The implementation checks rank-one support exactly; it does not replace
formal harmonic variables by numerical evaluations. Its deliberately
altered examples add a single pole of order $d+2$. This preserves the
nonzero critical coefficient while breaking the divisibility condition,
and the altered pencil is rejected in every recorded case.

\subsection{Generating the Stieltjes rows}
To obtain row $k$, only finite derivatives of $U_r,V_r,c_d$ through
orders $k$, $k+1$ are needed. Use \eqref{eq:cd-bell} for $c_d$,
ordinary finite differentiation for the endpoint powers, and the
coefficient rule \eqref{eq:Tjets-closed}. The next coefficient $b_{k+1}$
is compulsory because it multiplies the pole $u^{-1}$.

The exact suite tests this rule with independent formal symbols for the
Stieltjes constants. This verifies the coefficient extraction separately
from special-function numerics. Its analytic realization as an ordinary
series is the distinct content of Lemma~\ref{lem:block} and
Theorem~\ref{thm:subtraction}.

\subsection{Numerical representations used for independent checks}
For $0\le k\le d$, the moment integral is
\begin{equation}\label{eq:mellin-check}
 \mu_k(s,a)=\frac1{d!\Gamma(s)}\int_0^\infty
       t^{s-1}e^{-at}(1-e^{-t})^d
       \left(\sum_{n\ge0}n^k e^{-nt}\right)\dd t.
\end{equation}
For $k=0$, the series in parentheses is $(1-e^{-t})^{-1}$;
for $k>0$, it is a rational function obtained by repeated application
of $q\partial_q$ to $(1-q)^{-1}$ at $q=e^{-t}$. Thus this evaluator
uses only an elementary Mellin kernel and Gamma normalization, not the
finite endpoint formula being tested.

The numerical program computes the determinant of these integral moments
and compares it with $U_r+V_r\zeta$. Near $t=0$, it evaluates
$1-e^{-t}$ through an exponential-minus-one routine to avoid cancellation.
The hard-cutoff diagnostics instead use the finite endpoint formula with
increased precision. Their errors include actual finite-cutoff remainders,
so their reported accuracy is intentionally different from that of the
Mellin comparisons.

\subsection{Rebuilding and verifying the delivery}
From the package directory, run
\begin{verbatim}
python verify_manifest.py
python verification/verify_exact.py
python verification/verify_numeric.py
python build.py
\end{verbatim}
The two mathematical verification commands regenerate their JSON result records. The build
script regenerates the standalone TeX and runs the LaTeX compiler. The
checksum verifier checks the delivered files before modification; after
regenerating timestamped result records or a PDF, expected byte hashes
can differ. No network access is needed for either verification suite.
The exact dependency versions and the TeX package requirements are
recorded in the accompanying files.

% END sections/11-algorithms.tex
% BEGIN references.tex
\begin{thebibliography}{99}
\addcontentsline{toc}{section}{References}

\bibitem{repo:main}
ProveIt Contributors,
\emph{Polylogarithms and their Arithmetic Bridges}, canonical manuscript,
main file and preface. Repository snapshot \pin.
\url{https://github.com/VladimirReshetnikov/ProveIt/tree/73f18ef013d5b74eafc883efb2c15210642df12a/Analysis/Polylogarithms/docs/manuscript}.

\bibitem{repo:dhr}
ProveIt Contributors,
\emph{Determinantal Harmonic Resonance: Gauss--Vandermonde finite parts,
Stieltjes jets, and depth-two polylogarithmic closure}, preserved report,
\texttt{article\_standalone.tex}, at the same snapshot.
\url{https://github.com/VladimirReshetnikov/ProveIt/tree/73f18ef013d5b74eafc883efb2c15210642df12a/Analysis/Polylogarithms/docs/reports/stieltjes-harmonic-resonance/continuations/determinantal-harmonic-resonance}.
The repeated-pole, harmonic-rank, and rational-grid questions in its
further-research section are the direct targets of this continuation.

\bibitem{repo:incoming}
ProveIt Contributors,
\emph{Incoming research packages}, directory catalogue and selected
status/readme material, inspected 11 October 2026 at \pin.
\url{https://github.com/VladimirReshetnikov/ProveIt/tree/73f18ef013d5b74eafc883efb2c15210642df12a/docs/incoming}.
The precise limits of interior inspection are stated in Section~\ref{sec:audit}.

\bibitem{repo:integration}
ProveIt Contributors,
\emph{Integration chapter},
canonical manuscript \texttt{chapters/07-integration.tex}, opening 280
source lines inspected at the same snapshot.

\bibitem{repo:discovery}
ProveIt Contributors,
\emph{Discovery and proof-status chapter}, canonical manuscript
\texttt{chapters/10-discovery.tex}, at the same snapshot.

\bibitem{chen-singer}
S. Chen and M. F. Singer,
\emph{Residues and Telescopers for Rational Functions},
arXiv:1203.4200 (2012).
\url{https://arxiv.org/abs/1203.4200}.

\bibitem{arreche-sitaula}
C. E. Arreche and H. P. Sitaula,
\emph{A computational approach to rational summability and its applications
via discrete residues}, arXiv:2503.15636 (2025).
\url{https://arxiv.org/abs/2503.15636}.
Cited for the broader discrete-residue and rational-summability framework,
not as a source of the determinant identities proved here.

\bibitem{dlmf:hurwitz}
NIST Digital Library of Mathematical Functions,
\S25.11, \emph{Hurwitz Zeta Function}; in particular the defining series,
shift and parameter-derivative formulas, polygamma values, and Lerch's
identity at spectral order zero.
\url{https://dlmf.nist.gov/25.11}.

\bibitem{dlmf:gamma}
NIST Digital Library of Mathematical Functions,
Chapter 5, \emph{Gamma Function}, especially \S5.2 and \S5.15.
\url{https://dlmf.nist.gov/5.2}; \url{https://dlmf.nist.gov/5.15}.

\bibitem{dlmf:gauss}
NIST Digital Library of Mathematical Functions,
\S15.4, equation 15.4.20, Gauss hypergeometric summation.
\url{https://dlmf.nist.gov/15.4.E20}.

\bibitem{dlmf:polylog}
NIST Digital Library of Mathematical Functions,
\S25.12, \emph{Polylogarithms}, and \S25.14, \emph{Lerch's Transcendent}.
\url{https://dlmf.nist.gov/25.12}; \url{https://dlmf.nist.gov/25.14}.
The multiple-polylogarithm convention is defined explicitly in
\eqref{eq:MPL}.

\end{thebibliography}

% END references.tex
\end{document}
